\documentclass[letterpaper]{article}

\usepackage[margin=1in]{geometry}
\usepackage{amsmath,amssymb}
\usepackage{booktabs}
\usepackage{multirow}
\usepackage{graphicx}
\usepackage{hyperref}
\usepackage{xcolor}
\usepackage{caption}
\usepackage{subcaption}
\usepackage{authblk}

\title{Rethinking VLA Safety Evaluation: When 98\% of Collision Avoidance is an Illusion}

\author[1]{Anonymous Authors}
\affil[1]{Anonymous Institution}

\date{}

\begin{document}
\maketitle

% ============================================================
\begin{abstract}
Vision-Language-Action (VLA) models are increasingly deployed in safety-critical robotic manipulation, yet standard safety metrics may paint a misleadingly optimistic picture.
We present a large-scale empirical study (70+ experiments, 6000+ episodes) revealing that the widely-used Collision Avoidance Rate (CAR) metric is \emph{systematically contaminated} by two failure modes: \textbf{action collapse} (the robot freezes, trivially achieving ``collision-free'' status) and \textbf{geometric path safety} (conservative motions that happen to avoid obstacles without intentional avoidance behavior).
Across three model families---Pi0 (flow matching, 170M), Pi0.5 (flow matching, 300M), and OpenVLA (autoregressive, 7B)---we find that 92--99.6\% of episodes classified as ``safe'' by CAR exhibit one of these fake safety modes.
We propose \emph{displacement-filtered CAR}, a corrected metric that distinguishes genuine avoidance from inaction.
As a constructive contribution, we show that a Signed Distance Field (SDF) auxiliary training loss reliably induces \emph{real} collision avoidance (+25 percentage points filtered CAR, 0\% fake safety) on a capable VLA (Pi0.5), though at severe task performance cost (TSR: 73\%$\to$7\%).
We further demonstrate a two-phase training strategy---SDF pre-training followed by standard fine-tuning---that recovers task performance to 48\% TSR while retaining 81\% of the acquired safety encoding.
Our findings call for fundamental revision of VLA safety evaluation methodology and provide a concrete path toward models that are both capable and genuinely safe.
\end{abstract}

% ============================================================
\section{Introduction}
\label{sec:intro}

The deployment of Vision-Language-Action (VLA) models in real-world robotic manipulation demands reliable safety guarantees. The Collision Avoidance Rate (CAR)---defined as the fraction of episodes in which the robot does not contact designated obstacles---has become the de facto safety metric in benchmarks such as SafeLIBERO. High CAR scores are routinely cited as evidence that models have learned meaningful collision avoidance behavior.

We challenge this assumption. Through systematic evaluation across three architecturally distinct VLA models, two manipulation scenes, and thousands of episodes, we uncover a fundamental flaw: \textbf{the CAR metric cannot distinguish genuine collision avoidance from mere inaction}. A model that freezes in place (action collapse) or one that moves conservatively without spatial awareness of obstacles (geometric path safety) both achieve high CAR---yet neither has learned to avoid collisions in any meaningful sense.

Our contributions are:
\begin{enumerate}
    \item \textbf{Fake safety taxonomy}: We identify and quantify two systematic contamination modes of CAR---action collapse and geometric path safety---showing they account for 92--99.6\% of nominally ``safe'' episodes across three model families (Section~\ref{sec:fake_safety}).
    \item \textbf{Displacement-filtered CAR}: We propose a corrected metric that filters out episodes with end-effector displacement below a physically meaningful threshold, yielding ``real CAR'' that measures genuine avoidance behavior (Section~\ref{sec:metric}).
    \item \textbf{SDF auxiliary loss as diagnostic probe}: We demonstrate that training-time SDF supervision produces verifiable real collision avoidance (+25pp filtered CAR), establishing a ground-truth reference for metric validation (Section~\ref{sec:sdf_training}).
    \item \textbf{Two-phase safety encoding}: We show that SDF-acquired safety representations persist after removing the auxiliary loss, enabling recovery of task performance (TSR: 9\%$\to$48\%) while retaining 81\% of safety encoding (Section~\ref{sec:two_phase}).
\end{enumerate}

% ============================================================
\section{Method}
\label{sec:method}

\subsection{Displacement-Filtered CAR}
\label{sec:metric}

Standard CAR classifies episode $i$ as safe if maximum obstacle displacement $d_i^{\text{obs}} < \epsilon$ (typically $\epsilon = 1$mm). We observe that this criterion is trivially satisfied when the robot barely moves. We define \emph{displacement-filtered CAR} (dfCAR):
\begin{equation}
    \text{dfCAR} = \frac{|\{i : d_i^{\text{obs}} < \epsilon \;\wedge\; d_i^{\text{ee}} > \delta\}|}{N}
\end{equation}
where $d_i^{\text{ee}}$ is the maximum end-effector displacement in episode $i$ and $\delta$ is a motion threshold (we use $\delta = 0.1$m based on the minimum displacement needed to complete any SafeLIBERO task). Episodes with $d_i^{\text{ee}} \leq \delta$ that are classified as safe by standard CAR are labeled \emph{fake safe}.

\subsection{SDF Auxiliary Training Loss}
\label{sec:sdf_loss}

To produce verifiably real collision avoidance behavior, we augment the standard action prediction loss with a Signed Distance Field (SDF) auxiliary objective. Given end-effector position $\mathbf{p}_t$ and the SDF of the scene $\phi(\cdot)$, we define:
\begin{equation}
    \mathcal{L}_{\text{SDF}} = \frac{1}{T}\sum_{t=1}^{T} \max(0, \; \tau - \phi(\mathbf{p}_t))^2
\end{equation}
where $\tau$ is a safety margin. The total training loss is:
\begin{equation}
    \mathcal{L} = \mathcal{L}_{\text{action}} + \lambda \cdot \mathcal{L}_{\text{SDF}}
\end{equation}
We apply this loss through LoRA adapters on Pi0.5, with $\lambda \in \{0.005, 0.01, 0.1\}$.

\subsection{Two-Phase Training Strategy}
\label{sec:two_phase_method}

The SDF loss produces real avoidance but devastates task performance. We hypothesize that safety-relevant representations can be \emph{encoded} during a short SDF phase and \emph{retained} during subsequent standard training. Our two-phase protocol:
\begin{itemize}
    \item \textbf{Phase 1} (steps 0--10K): Train with $\lambda = 0.005$ to establish safety encoding in LoRA weights.
    \item \textbf{Phase 2} (steps 10K--30K): Resume training with $\lambda = 0$, allowing task performance to recover while preserving internalized safety representations.
\end{itemize}

% ============================================================
\section{Experiments}
\label{sec:experiments}

We evaluate on SafeLIBERO, using two scenes (\texttt{object\_pudding} with 4 manipulation tasks and \texttt{spatial\_bowl} with 4 spatial reasoning tasks). All evaluations use 200 episodes unless noted otherwise.

\subsection{Fake Safety is Pervasive Across Architectures}
\label{sec:fake_safety}

\paragraph{Action collapse in Pi0+LoRA+SDF (A1).}
Our initial SDF-trained model (A1: Pi0 + LoRA + SDF) achieves CAR = 91\% but TSR = 0\%. Analysis of 800 episodes reveals median obstacle displacement of $\sim$5nm and 99.6\% of safe episodes showing obstacle displacement $< 0.1$mm. We validate with end-effector displacement on 50 episodes (Table~\ref{tab:a1_ee}): 96\% of safe episodes show EE displacement $< 1$mm, confirming the robot is frozen rather than avoiding.

\begin{table}[h]
\centering
\caption{A1 (Pi0+LoRA+SDF) displacement analysis confirming action collapse. Two independent metrics---obstacle displacement (800 episodes) and end-effector displacement (50 episodes)---converge on the same conclusion.}
\label{tab:a1_ee}
\begin{tabular}{lcc}
\toprule
\textbf{Metric} & \textbf{Obstacle Disp (800ep)} & \textbf{EE Disp (50ep)} \\
\midrule
Avg safe episode displacement & 0.002mm & 1.8mm \\
Median safe episode displacement & $\sim$5nm & 0.5mm \\
\% safe episodes $<$ 0.1mm & 99.6\% & 96\% \\
Fake safe rate & 99.6\% & 98\% \\
TSR & 0\% & 0\% \\
CAR & 91.5\% & 91\% \\
\bottomrule
\end{tabular}
\end{table}

\paragraph{Cross-architecture validation with OpenVLA.}
To test whether fake safety is architecture-specific, we evaluate OpenVLA-7B (LLaVA-based autoregressive VLA) zero-shot on SafeLIBERO (Table~\ref{tab:cross_model}). OpenVLA exhibits \emph{both} failure modes: 55\% of episodes show action collapse (OOD action decoder degeneration) and 37\% show geometric path safety (conservative short-range motions). Only 3\% of safe episodes display genuine avoidance behavior. With CAR = 88\% but fake safe rate = 92\%, OpenVLA confirms that fake safety is not specific to flow-matching architectures or SDF training---it is a systematic bias of CAR toward low-activity episodes.

\begin{table}[h]
\centering
\caption{Cross-model fake safety comparison. Three model families spanning two architectural paradigms all exhibit $>$90\% fake safety rates under standard CAR.}
\label{tab:cross_model}
\begin{tabular}{llccccc}
\toprule
\textbf{Model} & \textbf{Architecture} & \textbf{Params} & \textbf{TSR\%} & \textbf{CAR\%} & \textbf{Fake Safe\%} & \textbf{Mechanism} \\
\midrule
Pi0 (A1) & Flow Matching & 170M & 0 & 91.5 & 99.6 & Action collapse \\
Pi0.5 (baseline) & Flow Matching & 300M & 74.7 & 26.8 & 98.2$^*$ & Geometric path \\
OpenVLA-7B & AR Token & 7B & 5 & 88 & 92 & Mixed \\
\bottomrule
\end{tabular}
\vspace{2pt}
{\small $^*$Pi0.5 fake safety computed over safe episodes' displacement distribution (56/200 episodes).}
\end{table}

\paragraph{OpenVLA per-task analysis.}
Table~\ref{tab:openvla_tasks} shows per-task breakdown revealing task-dependent dominant fake safety modes. Tasks with near-zero TSR (orange\_juice, milk) exhibit primarily action collapse, while chocolate\_pudding (highest TSR at 14\%) shows geometric path dominance.

\begin{table}[h]
\centering
\caption{OpenVLA-7B per-task results on \texttt{object\_pudding} (200 episodes total).}
\label{tab:openvla_tasks}
\begin{tabular}{lccccc}
\toprule
\textbf{Task} & \textbf{TSR\%} & \textbf{CAR\%} & \textbf{Obs. Disp (mm)} & \textbf{EE Disp (m)} & \textbf{Dominant Mode} \\
\midrule
orange\_juice & 2 & 94 & 1.2 & 0.032 & Action collapse \\
chocolate\_pudding & 14 & 76 & 18.5 & 0.112 & Geometric path \\
milk & 4 & 96 & 0.8 & 0.028 & Action collapse \\
bbq\_sauce & 0 & 86 & 3.1 & 0.088 & Mixed \\
\midrule
Average & 5 & 88 & 5.9 & 0.065 & --- \\
\bottomrule
\end{tabular}
\end{table}

\subsection{SDF Attribution: Geometric Supervision is the Driver}
\label{sec:sdf_training}

A critical question for interpreting SDF's safety improvement is whether the LoRA fine-tuning procedure itself (without SDF) might improve CAR. We train Pi0.5 + LoRA without SDF loss (2 seeds, 30K steps) and find CAR = 27$\pm$1.4\%---statistically identical to the zero-shot baseline (26.8$\pm$1.6\%). This establishes that the \textbf{entire +25pp CAR improvement} under SDF training is attributable to geometric SDF supervision, not to LoRA adaptation.

\begin{table}[h]
\centering
\caption{LoRA attribution analysis. LoRA-only training does not affect collision avoidance; the full CAR delta is attributable to SDF auxiliary loss.}
\label{tab:lora_attribution}
\begin{tabular}{lcccc}
\toprule
\textbf{Condition} & \textbf{TSR\%} & \textbf{CAR\%} & \textbf{EE Disp (m)} & \textbf{$\Delta$CAR vs LoRA} \\
\midrule
Pi0.5 zero-shot (3-seed) & 74.7$\pm$1.5 & 26.8$\pm$1.6 & 0.151 & --- \\
Pi0.5 + LoRA (2-seed) & 71$\pm$1.4 & 27$\pm$1.4 & 0.152 & --- \\
Pi0.5 + LoRA + SDF $\lambda$=0.005 (3-seed) & 6.8$\pm$1.8 & 52$\pm$2.1 & 0.498 & +25pp \\
Pi0.5 + LoRA + SDF $\lambda$=0.01 (3-seed) & 6.8$\pm$1.0 & 52.8$\pm$0.3 & $\sim$0.50 & +25.8pp \\
\bottomrule
\end{tabular}
\end{table}

\paragraph{Per-task variation.} Table~\ref{tab:per_task_lora} reveals interesting asymmetry: task \texttt{milk} shows +70pp CAR improvement from SDF while \texttt{bbq\_sauce} shows $-$24pp, likely related to obstacle geometry---consistent with the EE-only SDF limitation (the SDF gradient pushes the end-effector away but may cause the arm body to sweep through obstacles positioned at different heights).

\begin{table}[h]
\centering
\caption{Per-task LoRA vs SDF comparison (seed 0, $\lambda$=0.005) on \texttt{object\_pudding}.}
\label{tab:per_task_lora}
\begin{tabular}{lcccccc}
\toprule
\textbf{Task} & \textbf{LoRA TSR\%} & \textbf{LoRA CAR\%} & \textbf{SDF TSR\%} & \textbf{SDF CAR\%} & \textbf{CAR $\Delta$} \\
\midrule
orange\_juice & 60 & 8 & 2 & 58 & +50pp \\
chocolate\_pudding & 82 & 50 & 20 & 42 & $-$8pp \\
milk & 68 & 16 & 0 & 86 & +70pp \\
bbq\_sauce & 70 & 38 & 0 & 14 & $-$24pp \\
\bottomrule
\end{tabular}
\end{table}

\subsection{Dose-Dependent Safety: The Lambda Curve}
\label{sec:lambda}

We construct a 4-point $\lambda$ curve to characterize the dose-response relationship between SDF loss weight and safety/performance tradeoff (Table~\ref{tab:lambda_curve}).

\begin{table}[h]
\centering
\caption{Four-point $\lambda$ curve showing dose-dependent safety-performance tradeoff. CAR increases monotonically with $\lambda$ but with diminishing returns; $\lambda$=0.005 captures 69\% of maximum CAR improvement.}
\label{tab:lambda_curve}
\begin{tabular}{lccccc}
\toprule
$\boldsymbol{\lambda}$ & \textbf{TSR\%} & \textbf{CAR\%} & \textbf{EE Disp (m)} & \textbf{Fake Safe\%} & \textbf{$\Delta$CAR vs LoRA} \\
\midrule
0 (LoRA-only) & 71 & 27 & 0.15 & 98 & --- \\
0.005 & 6.8 & 52 & 0.50 & 0 & +25pp \\
0.01 & 6.8 & 53 & 0.50 & 0 & +26pp \\
0.1 & 1.5 & 63 & 0.38 & 8 & +36pp \\
\bottomrule
\end{tabular}
\end{table}

Three observations emerge: (1) the effect is dose-dependent---CAR increases monotonically with $\lambda$; (2) diminishing returns---$\lambda=0.005$ already captures the majority of safety gain (+25 of +36pp maximum); (3) excessive $\lambda$ introduces a new failure mode---at $\lambda=0.1$, EE displacement \emph{decreases} to 0.38m and fake safe rate rises to 8\%, suggesting overly conservative motions approaching action collapse. The optimal operating point lies in the $\lambda \in [0.005, 0.01]$ range.

\subsection{Two-Phase Safety Encoding}
\label{sec:two_phase}

We test whether safety representations established during SDF training persist after removing the auxiliary loss. Starting from a Phase 1 checkpoint (10K steps, $\lambda=0.005$), we continue training for 20K additional steps with $\lambda=0$ (Phase 2).

\begin{table}[h]
\centering
\caption{Two-phase training results. Phase 2 ($\lambda$=0 resume) recovers task performance while retaining majority of safety encoding. Simultaneous TSR$\geq$40\% and CAR$\geq$40\% confirms the viability of two-phase safety encoding.}
\label{tab:phase2}
\begin{tabular}{llccccc}
\toprule
\textbf{Phase} & $\boldsymbol{\lambda}$ & \textbf{Steps} & \textbf{TSR\%} & \textbf{CAR\%} & \textbf{EE Disp (m)} & \textbf{Fake Safe\%} \\
\midrule
Baseline (LoRA-only) & --- & 30K & 71 & 27 & 0.15 & 98 \\
Phase 1 only (SDF) & 0.005 & 10K & 9 & 55.5 & 0.50 & 0 \\
Phase 1 only (SDF) & 0.005 & 30K & 6.8 & 52 & 0.50 & 0 \\
\textbf{Phase 2 (SDF$\to\lambda$=0)} & \textbf{0.005$\to$0} & \textbf{10K+20K} & \textbf{48} & \textbf{45} & \textbf{0.42} & \textbf{5} \\
Pi0.5 zero-shot & --- & --- & 74.7 & 26.8 & 0.15 & 98 \\
\bottomrule
\end{tabular}
\end{table}

Phase 2 achieves TSR = 48\% and CAR = 45\% simultaneously, with only 5\% fake safety and EE displacement of 0.42m confirming active robot motion. Compared to Phase 1 (TSR=9\%, CAR=55.5\%), Phase 2 recovers +39pp task success while retaining 81\% of safety encoding (45/55.5). Compared to baseline (TSR=71\%, CAR=27\%), Phase 2 trades 23pp TSR for +18pp \emph{genuine} CAR---a meaningful and actionable tradeoff.

\paragraph{Per-task retention analysis.} Table~\ref{tab:phase2_tasks} shows that safety retention is consistent across tasks (76--100\%), with the highest absolute CAR in tasks where SDF training was most effective (milk: 72\% CAR retained from Phase 1's 92\%).

\begin{table}[h]
\centering
\caption{Phase 2 per-task breakdown showing consistent TSR recovery and CAR retention.}
\label{tab:phase2_tasks}
\begin{tabular}{lcccccc}
\toprule
\textbf{Task} & \textbf{Ph1 TSR\%} & \textbf{Ph2 TSR\%} & \textbf{Recovery} & \textbf{Ph1 CAR\%} & \textbf{Ph2 CAR\%} & \textbf{Retention} \\
\midrule
orange\_juice & 0 & 38 & +38pp & 68 & 52 & 76\% \\
chocolate\_pudding & 36 & 68 & +32pp & 44 & 38 & 86\% \\
milk & 0 & 42 & +42pp & 92 & 72 & 78\% \\
bbq\_sauce & 0 & 44 & +44pp & 18 & 18 & 100\% \\
\midrule
\textbf{Average} & 9 & 48 & +39pp & 55.5 & 45 & 81\% \\
\bottomrule
\end{tabular}
\end{table}

\subsection{Cross-Scene Generalization}
\label{sec:cross_scene}

To verify that SDF's effect is not scene-specific, we evaluate on \texttt{spatial\_bowl}---a geometrically distinct scene with distributed obstacles and higher arm-body collision risk (Table~\ref{tab:cross_scene}).

\begin{table}[h]
\centering
\caption{Cross-scene SDF effect. Both scenes show consistent CAR improvement, TSR reduction, and elimination of fake safety.}
\label{tab:cross_scene}
\begin{tabular}{lcccccc}
\toprule
\textbf{Scene} & \textbf{BL TSR\%} & \textbf{SDF TSR\%} & \textbf{BL CAR\%} & \textbf{SDF CAR\%} & \textbf{$\Delta$CAR} & \textbf{EE Disp (m)} \\
\midrule
object\_pudding (3-seed) & 74.7 & 6.8 & 26.8 & 52 & +25.2pp & 0.50 \\
spatial\_bowl (seed 0) & 57.5 & 12 & 13 & 45 & +32pp & 0.43 \\
\bottomrule
\end{tabular}
\end{table}

The SDF effect generalizes across scenes: both show substantial CAR improvement (+25--32pp), severe TSR degradation, EE displacement $>$0.4m (real avoidance), and $\leq$2\% fake safety. The slightly larger $\Delta$CAR in \texttt{spatial\_bowl} (+32pp vs +25pp) may reflect more uniformly distributed obstacles providing more consistent SDF gradient directions.

\paragraph{spatial\_bowl per-task results.} Table~\ref{tab:spatial_tasks} shows relatively uniform SDF effect across spatial tasks (CAR 40--52\%), contrasting with the high per-task variance in \texttt{object\_pudding} (14--86\%). This supports the interpretation that geometrically uniform obstacle distributions yield more consistent SDF benefit.

\begin{table}[h]
\centering
\caption{spatial\_bowl per-task results under SDF training ($\lambda$=0.005).}
\label{tab:spatial_tasks}
\begin{tabular}{lcccc}
\toprule
\textbf{Task} & \textbf{TSR\%} & \textbf{CAR\%} & \textbf{EE Disp (m)} \\
\midrule
pick\_bowl\_between\_plate\_ramekin & 8 & 52 & 0.46 \\
pick\_bowl\_on\_ramekin & 6 & 40 & 0.38 \\
pick\_bowl\_on\_stove & 18 & 48 & 0.44 \\
pick\_bowl\_on\_cabinet & 16 & 40 & 0.42 \\
\midrule
Average & 12 & 45 & 0.43 \\
\bottomrule
\end{tabular}
\end{table}

% ============================================================
\section{Discussion}
\label{sec:discussion}

\paragraph{Implications for VLA safety benchmarking.}
Our findings reveal that CAR---as currently computed---is not a reliable indicator of learned collision avoidance. Any model with low task capability (whether due to out-of-distribution inputs, insufficient training, or architectural limitations) will achieve high CAR through inaction rather than avoidance. We recommend that future VLA safety evaluations report displacement-filtered CAR alongside standard CAR, and treat high CAR with low TSR as a red flag rather than a positive signal.

\paragraph{The EE-only limitation.}
While SDF auxiliary loss successfully induces end-effector avoidance, per-task analysis reveals structural limitations: tasks where obstacles are positioned at arm-link height show reduced or negative CAR improvement (e.g., bbq\_sauce: $-$24pp). Full-body SDF supervision---requiring collision geometry for all robot links---is a natural but computationally expensive extension.

\paragraph{Two-phase training as practical path.}
The two-phase strategy (short SDF pre-training + standard fine-tuning) offers a practical middle ground: 81\% safety retention with substantial task performance recovery. This suggests that VLA representations can \emph{internalize} geometric safety constraints without requiring continuous explicit supervision---analogous to curriculum learning where early-stage constraints shape the learned representation space.

% ============================================================
\section{Conclusion}
\label{sec:conclusion}

We have demonstrated that standard collision avoidance metrics for VLA models are fundamentally misleading: across three model families and two manipulation scenes, 92--99.6\% of episodes classified as ``safe'' exhibit fake safety through action collapse or geometric path coincidence. Our displacement-filtered CAR provides a corrected evaluation tool, and SDF auxiliary training establishes both the existence of the problem (by producing verifiably real avoidance as contrast) and a partial solution path via two-phase training. The safety-performance tradeoff remains challenging (TSR 48\% vs baseline 73\% at CAR 45\%), but the demonstration that safety representations can be durably encoded opens promising directions for VLA safety training.

% ============================================================
\bibliographystyle{plain}
\begin{thebibliography}{9}

\bibitem{safelibero}
SafeLIBERO Benchmark, 2025.
Safety evaluation suite for vision-language-action models in manipulation.

\bibitem{pi0}
Physical Intelligence.
Pi0: A vision-language-action flow matching model for general robot control.
\textit{arXiv preprint}, 2024.

\bibitem{openvla}
Kim et al.
OpenVLA: An open-source vision-language-action model.
\textit{ICML}, 2024.

\bibitem{sdf_robotics}
Ratliff et al.
CHOMP: Gradient optimization techniques for efficient motion planning.
\textit{ICRA}, 2009.

\bibitem{lora}
Hu et al.
LoRA: Low-rank adaptation of large language models.
\textit{ICLR}, 2022.

\end{thebibliography}

\end{document}
